\providecommand{\Mtf}{M_{24}}
\providecommand{\cent}{C_{\Mtf}}
\providecommand{\Res}{\mathrm{Res}}
\providecommand{\Bor}{\operatorname{Bor}}
\providecommand{\wt}{\operatorname{wt}}
\providecommand{\cO}{\mathcal O}
\providecommand{\fg}{\mathfrak g}
\providecommand{\ZZ}{\mathbb Z}
\providecommand{\kBKM}{\kappa_{\mathrm{BKM}}}
\providecommand{\kch}{\kappa_{\mathrm{ch}}}
\providecommand{\kcat}{\kappa_{\mathrm{cat}}}
\providecommand{\kfib}{\kappa_{\mathrm{fiber}}}
\providecommand{\kHeis}{\kappa_{\mathrm{ch}}^{\mathrm{Heis}}}

\section*{Order-two Mathieu centralisers and the CHL weight}
\label{note:sec:n2-2a-2b-anomalies}

\paragraph{The two involution classes.}
The order-two Mathieu classes carry different data.  In ATLAS notation
\[
\begin{array}{c|c|c|c}
\text{class} & \text{frame shape} & \cent(g) &
H^3(\cent(g);U(1))\\
\hline
2A & 1^8 2^8 & 2^{1+4}_{+}{:}(S_3\times S_3) & \ZZ/2\\
2B & 2^{12} & 2^4{:}A_8 & \ZZ/4 .
\end{array}
\]
The first line is the Mukai-Nikulin symplectic involution used in the
primitive level-two CHL denominator.  The second line is the fixed-point
free Mathieu involution on the \(24\)-point permutation representation.
Milgram's prime-local computation of \(H^\ast(M_{24};\ZZ)\), together
with Benson's Sylow restriction theorem, gives the two residual
centraliser groups displayed above.  Let
\(\xi_{2A}\) and \(\eta_{2B}\) denote fixed generators of these cyclic
groups.

The cyclic Mathieu anomaly of Gaberdiel-Persson-Volpato restricts
trivially to \(\langle 2A\rangle\) and \(\langle 2B\rangle\):
\[
 \Res^{M_{24}}_{\langle g\rangle}[\alpha_{K3}]=0,\qquad
 g\in\{2A,2B\}.
\]
Thus the primitive cyclic CHL construction at \(N=2\) is anomaly-free.
The classes \(\xi_{2A}\) and \(\eta_{2B}\) are centraliser-level
discrete torsion for twisted-twined sectors.  Under the
Dijkgraaf-Witten pointed-fusion dictionary they define associator twists
\[
 \mathrm{Vec}_{\cent(2A)}^{\xi_{2A}},\qquad
 \mathrm{Vec}_{\cent(2B)}^{\eta_{2B}},
\]
and their Drinfeld centres are twisted quantum doubles of the two
centralisers.  This changes the multiplier system of a twisted-twined
partition function.  It leaves the Borcherds weight projection untouched.

\paragraph{Ordinary Mathieu twining.}
For a Mathieu class \([g]\subset M_{24}\), the ordinary twined K3
elliptic genus is a weak Jacobi form
\[
 Z^{(g)}_{K3}(\tau,z)
  =\operatorname{Tr}_{\mathcal H_{RR}}
  \left(g(-1)^{F+\widetilde F}q^{L_0-c/24}
  \bar q^{\widetilde L_0-\widetilde c/24}y^{J_0}\right)
  \in J^{w,\nu_g}_{0,1}(\Gamma_0(|g|)).
\]
In the Gaberdiel-Hohenegger-Volpato normalisation
\[
 Z^{(g)}_{K3}(\tau,z)=2(y^{-1}+y)+c^{(g)}(0,0)+O(q),
\]
the two order-two coefficients are
\[
\begin{array}{c|cc}
g & 2A & 2B\\
\hline
c^{(g)}(0,0) & 4 & -4\\
c^{(g)}(0,0)/2 & 2 & -2 .
\end{array}
\]
These are ordinary single-twining coefficients, separate from the
primitive CHL constants.  The \(2B\) entry has negative single-twining
scalar weight, and the Gritsenko lift stops before any positive
holomorphic primitive BKM denominator.  The primitive level-two CHL row
uses \(2A\), together with the CHL cusp normalisation below.

\paragraph{Primitive CHL normalisation.}
Let \(g_2\) be the class \(2A\), with frame shape \(1^82^8\).  In the
level-two CHL denominator normalisation the weak Jacobi input
\(\varphi^{\mathrm{CHL}}_2\) has
\[
 \varphi^{\mathrm{CHL}}_2(\tau,z)\big|_{q^0}
  = \zeta^{-1}+8+\zeta,\qquad \zeta=e^{2\pi iz}.
\]
Equivalently
\[
 c_2(0):=[q^0\zeta^0]\varphi^{\mathrm{CHL}}_2=8.
\]
The comparison map from the Mathieu class to the CHL denominator is
\[
(2A,\ 1^82^8)
 \xrightarrow{\mathrm{GHV}}
 Z^{(2A)}_{K3}
 \xrightarrow{\mathrm{CHL\ cusp}}
 \varphi^{\mathrm{CHL}}_2
 \xrightarrow{\Bor}
 \Phi_2=\Delta^{(2)}_4
 \xrightarrow{\wt}
 \kBKM(\Phi_2)=4 .
\]
The coefficient calculation is
\[
 T_{H^\ast}(g_2)=16,\qquad A_2=4,\qquad
 c_2(0)=T_{H^\ast}(g_2)-2A_2=16-8=8.
\]
Here \(T_{H^\ast}(g_2)\) is the Mukai-lattice Lefschetz trace of the
frame shape \(1^82^8\), and \(A_2\) is the odd theta-component
coefficient in the Eichler-Zagier decomposition.

Borcherds' singular theta lift gives
\[
 \wt(\Bor(\varphi))=\frac{c(0)}{2}.
\]
Therefore
\[
 \kBKM(\Phi_2)=\wt(\Phi_2)=\frac{c_2(0)}2=4 .
\]

\paragraph{The primitive CHL row.}
For the primitive CHL branch
\[
 N\in\{1,2,3,4,6\},
\]
the Borcherds input constants and weights are
\[
\begin{array}{c|ccccc}
N & 1 & 2 & 3 & 4 & 6\\
\hline
c_N(0) & 10 & 8 & 6 & 4 & 2\\
\kBKM(\Phi_N)=c_N(0)/2 & 5 & 4 & 3 & 2 & 1 .
\end{array}
\]
This table is the primitive CHL row.  Lower twining, row-normalised, or
non-primitive constants require an explicit row map before comparison.

\paragraph{The level-one Igusa square.}
At level one the Jacobi input is
\[
 \phi_{0,1}(\tau,z)=\zeta^{-1}+10+\zeta+O(q).
\]
The Borcherds product is
\[
 \Delta_5=\Bor(\phi_{0,1}),\qquad c_1(0)=10,\qquad
 \kBKM(\Delta_5)=5.
\]
The monic denominator is
\[
 D_5=64^{-1}\Delta_5,\qquad
 \operatorname{den}(\fg_{\Delta_5})=D_5(2Z).
\]
The square normalisation of the Igusa form is
\[
 \Phi_{10}^{\mathrm{sq}}=\Delta_5^2,
\]
while the Oberdieck-Pandharipande scalar convention uses
\[
 \Phi_{10}^{\mathrm{OP}}=D_5^2=4096^{-1}\Delta_5^2,\qquad
 Z_{\mathrm{OP/DT}}=-(\Phi_{10}^{\mathrm{OP}})^{-1}
  =-4096\,\Delta_5^{-2}.
\]
The scalar \(64\) fixes the leading theta coefficient of the monic
product.  It leaves the weight \(\kBKM(\Delta_5)=5\) unchanged.

\paragraph{Invariant separation on \(K3\times E\).}
For \(Y=K3\times E\),
\[
 \kcat(Y)=\chi(\cO_{K3})\chi(\cO_E)=2\cdot0=0,
\]
and the compact chiral Hodge supertrace is
\[
 \kch(Y)=\sum_q(-1)^q h^{0,q}(Y)=0.
\]
The Heisenberg specialisation, Borcherds denominator, and K3 fibre
rank are distinct constructions:
\[
 \kHeis(Y)=3,\qquad
 \kBKM(\fg_{\Delta_5})=5,\qquad
 \kfib(Y)=24.
\]
These invariants live in separate constructions.  The formula for
\(\kBKM\) used here is the automorphic denominator weight \(c_N(0)/2\);
\(\kcat\) and \(\kch\) are total-space Hodge invariants.

\paragraph{Branch and operadic scope.}
The BKM branch for \(K3\times E\) is the Hall-Drinfeld or Borcherds
denominator branch, represented at level one by \(\fg_{\Delta_5}\) and
at level two by \(\fg_{\Phi_2}\).  The Yangian branch is the Mukai
self-mirror branch, governed by rank-\(24\) Mukai parameters and
Maulik-Okounkov stable envelopes.  The order-two comparison above is a
trace-to-denominator map through a Jacobi input.  It identifies the
centraliser multiplier contribution entering the CHL denominator
weight.  The Hall-Drinfeld or Borcherds branch and the Mukai Yangian
branch remain different objects.

Since \(K3\times E\) is a Calabi-Yau threefold, the native specialised
\(\Phi\)-output is \(E_1\)-chiral.  An \(E_2\)-structure belongs to the
centre, Drinfeld double, or module layer associated to the \(E_1\)
object.  The order-two centraliser torsion changes multiplier data at
that auxiliary layer and leaves the CHL weight identity
\[
 \kBKM(\Phi_2)=c_2(0)/2=4
\]
unchanged.

\paragraph{References.}
ATLAS centralisers for \(M_{24}\) fix the two order-two centraliser
groups.  Milgram \(1995\) supplies the prime-local cohomology
detection, and Benson \(1998\) supplies the Sylow restriction theorem.
Eguchi-Ooguri-Tachikawa, Gaberdiel-Hohenegger-Volpato, and
Eguchi-Hikami fix the Mathieu twining normalisation.  Borcherds'
singular theta theorem gives \(\wt=c(0)/2\).  Gritsenko and Nikulin
construct the CHL denominator ladder and the \(\Delta_5\) denominator.
Oberdieck-Pandharipande fix the reduced \(K3\times E\) scalar square.
